\documentclass[12pt]{article}

\usepackage{amsmath, amsthm, amssymb}
\usepackage{geometry}
\usepackage{hyperref}
\usepackage{booktabs}
\usepackage{array}

\geometry{margin=1in}

\newtheorem{theorem}{Theorem}[section]
\newtheorem{lemma}[theorem]{Lemma}
\newtheorem{corollary}[theorem]{Corollary}
\newtheorem{proposition}[theorem]{Proposition}
\newtheorem{definition}[theorem]{Definition}
\newtheorem{observation}[theorem]{Observation}

\title{Digital Root Algebra and the 37-Field:\\
A Unified Framework Connecting\\
Prime Structure, Pascal's Triangle, and Tesla FLUX States}

\author{Michael Warren Song\\
\small\texttt{red3rdeye@gmail.com}}

\date{2026}

\begin{document}

\maketitle

\begin{abstract}
We develop a unified framework built on the digital root function and a modular
field centered on 37, connecting four areas: (1) a provable structural constraint
on twin prime pairs via digital root residue classes; (2) the canonical mod-37
degeneracy of the digit blocks $\{123, 456, 789\}$; (3) the 191--137 prime bridge,
in which both primes share digital root 2 and each leads a twin prime pair with
the same digital root pattern $(2,4)$; and (4) a unified analysis of Pascal's
Row 8 demonstrating that the first three elements sum to exactly 37, that the
adjacent-pair-sum operation produces 100\% occupancy in the Tesla FLUX states
$\{3,6,9\}$, and that the product of element digital roots reduces through 19
to 1.  The framework reveals 12 --- equal to $123 \bmod 37$ --- as a fixed
point of the digit-sum operation, and identifies the Tesla FLUX class $\{3,6,9\}$
as the structurally excluded zone for twin primes greater than 3.
\end{abstract}

\tableofcontents
\newpage

%──────────────────────────────────────────────────────────────────────────────
\section{Introduction}
\label{sec:intro}
%──────────────────────────────────────────────────────────────────────────────

The digital root of a positive integer $n$, written $\mathrm{DR}(n)$, is the
single digit obtained by iterating the digit-sum map until a value in $\{1,\ldots,9\}$
is reached.  It is well known that $\mathrm{DR}(n) = 1 + (n-1)\bmod 9$ for
$n \geq 1$, making the digital root equivalent to the representative of $n$ in
$\mathbb{Z}/9\mathbb{Z}$ with the convention that 0 is replaced by 9.

Despite its elementary definition, the digital root carries non-trivial structural
information about primes, combinatorial sequences, and modular arithmetic.  This
paper assembles four interconnected results around a common axis: the mod-37
field and the Tesla FLUX states $\{3,6,9\}$.

\textbf{Organisation.}  Section~\ref{sec:prelim} fixes notation and definitions.
Section~\ref{sec:twin} proves the DR constraint on twin primes.
Section~\ref{sec:37field} establishes the canonical mod-37 degeneracy.
Section~\ref{sec:bridge} presents the 191--137 bridge observations.
Section~\ref{sec:pascal} analyses Pascal's Row 8 under all three lenses
simultaneously.  Section~\ref{sec:digitsum} proves the digit-sum invariant.
Section~\ref{sec:drproduct} traces the DR-product chain for Row~8.
Section~\ref{sec:conclusion} collects open questions.

%──────────────────────────────────────────────────────────────────────────────
\section{Preliminaries}
\label{sec:prelim}
%──────────────────────────────────────────────────────────────────────────────

\begin{definition}[Digital Root]
For $n \in \mathbb{Z}_{\geq 0}$, define
\[
    \mathrm{DR}(n) =
    \begin{cases}
        0 & \text{if } n = 0, \\
        1 + (n-1)\bmod 9 & \text{if } n \geq 1.
    \end{cases}
\]
\end{definition}

The digital root satisfies $\mathrm{DR}(a+b) \equiv \mathrm{DR}(a)+\mathrm{DR}(b)
\pmod{9}$ for all $a,b\geq 0$, where the result is again taken in $\{1,\ldots,9\}$
(replacing 0 by 9).  In particular:

\begin{lemma}[9 is the DR identity]
\label{lem:identity}
For all $n \geq 1$ and $k \geq 0$,\; $\mathrm{DR}(n + 9k) = \mathrm{DR}(n)$.
\end{lemma}

\begin{proof}
$\mathrm{DR}(n+9k) = 1 + (n+9k-1)\bmod 9 = 1 + (n-1)\bmod 9 = \mathrm{DR}(n)$.
\end{proof}

\begin{definition}[Tesla FLUX States]
Define the \emph{Tesla FLUX class} $\mathcal{F} = \{3, 6, 9\}$.
A positive integer $n$ is \emph{FLUX-active} if $\mathrm{DR}(n) \in \mathcal{F}$,
equivalently if $3 \mid n$.  The \emph{spine FLUX ratio} of a finite multiset
$S$ of positive integers is $|\{s \in S : \mathrm{DR}(s)\in\mathcal{F}\}|/|S|$.
\end{definition}

By Lemma~\ref{lem:identity}, the set $\mathcal{F}$ is closed under addition of
any multiple of 9: if $\mathrm{DR}(n)\in\mathcal{F}$ then $\mathrm{DR}(n+9k)
\in \mathcal{F}$ for all $k\geq 0$.  Moreover, $\mathcal{F}$ generates itself
under repeated addition of 3:
\[
    3 \xrightarrow{+3} 6 \xrightarrow{+3} 9 \xrightarrow{+3} 12 \xrightarrow{\mathrm{DR}} 3 \xrightarrow{\cdots}
\]

\begin{observation}
$1+2+3 = 6 \in \mathcal{F}$ and $1\cdot2\cdot3 = 6 \in \mathcal{F}$.
The sum and product of $\{1,2,3\}$ coincide and land in the Tesla FLUX class.
\end{observation}

\begin{definition}[37-Field]
We work in $\mathbb{Z}/37\mathbb{Z}$, referred to as the \emph{37-field}.
For a positive integer $n$ we write $[n]_{37} = n \bmod 37$.
\end{definition}

%──────────────────────────────────────────────────────────────────────────────
\section{DR Constraint on Twin Primes}
\label{sec:twin}
%──────────────────────────────────────────────────────────────────────────────

\begin{theorem}[DR Patterns of Twin Primes]
\label{thm:twin}
Let $(p, p+2)$ be a twin prime pair with $p > 3$.  Then
\[
    \bigl(\mathrm{DR}(p),\, \mathrm{DR}(p+2)\bigr) \;\in\;
    \{(2,4),\;(5,7),\;(8,1)\}.
\]
In particular, neither $p$ nor $p+2$ is FLUX-active.
The Tesla FLUX class $\mathcal{F} = \{3,6,9\}$ is the
\emph{excluded zone} for twin prime digital roots when $p > 3$.
\end{theorem}

\begin{proof}
Since $p > 3$ is prime, $p$ is not divisible by 3, so $\mathrm{DR}(p) \notin
\mathcal{F}$, i.e.\ $\mathrm{DR}(p) \in \{1,2,4,5,7,8\}$.  Similarly $p+2$
is prime and greater than 3, so $\mathrm{DR}(p+2) \notin \mathcal{F}$.

Since $\mathrm{DR}(p+2) \equiv \mathrm{DR}(p) + 2 \pmod{9}$ (with values in
$\{1,\ldots,9\}$), we have $\mathrm{DR}(p+2) = \mathrm{DR}(p) + 2$ if
$\mathrm{DR}(p) \leq 7$, and $\mathrm{DR}(p+2) = \mathrm{DR}(p) - 7$ if
$\mathrm{DR}(p) \geq 8$.

Checking each allowed value of $\mathrm{DR}(p)$:
\begin{center}
\begin{tabular}{ccc}
\toprule
$\mathrm{DR}(p)$ & $\mathrm{DR}(p+2)$ & Allowed? \\
\midrule
1 & 3 & \textbf{No} ($3 \in \mathcal{F}$, so $3\mid p+2$) \\
2 & 4 & Yes \\
4 & 6 & \textbf{No} ($6 \in \mathcal{F}$, so $3\mid p+2$) \\
5 & 7 & Yes \\
7 & 9 & \textbf{No} ($9 \in \mathcal{F}$, so $3\mid p+2$) \\
8 & 1 & Yes \\
\bottomrule
\end{tabular}
\end{center}
Exactly three patterns survive.
\end{proof}

\begin{corollary}
\label{cor:twin_dr2}
If $(p,p+2)$ is a twin prime pair with $p > 3$ and $\mathrm{DR}(p)=2$, then
$\mathrm{DR}(p+2)=4$.
\end{corollary}

Table~\ref{tab:twin} verifies Theorem~\ref{thm:twin} for the first twenty
twin prime pairs.

\begin{table}[h]
\centering
\begin{tabular}{rrccrr}
\toprule
$p$ & $p+2$ & $\mathrm{DR}(p)$ & $\mathrm{DR}(p+2)$ & Pattern & $p\bmod 37$ \\
\midrule
3   & 5   & 3 & 5 & (3,5) & 3 \\
5   & 7   & 5 & 7 & (5,7) & 5 \\
11  & 13  & 2 & 4 & (2,4) & 11 \\
17  & 19  & 8 & 1 & (8,1) & 17 \\
29  & 31  & 2 & 4 & (2,4) & 29 \\
41  & 43  & 5 & 7 & (5,7) & 4 \\
59  & 61  & 5 & 7 & (5,7) & 22 \\
71  & 73  & 8 & 1 & (8,1) & 34 \\
101 & 103 & 2 & 4 & (2,4) & 27 \\
107 & 109 & 8 & 1 & (8,1) & 33 \\
137 & 139 & 2 & 4 & (2,4) & 26 \\
149 & 151 & 5 & 7 & (5,7) & 1 \\
179 & 181 & 8 & 1 & (8,1) & 31 \\
191 & 193 & 2 & 4 & (2,4) & 6 \\
197 & 199 & 8 & 1 & (8,1) & 12 \\
227 & 229 & 2 & 4 & (2,4) & 5 \\
239 & 241 & 5 & 7 & (5,7) & 17 \\
269 & 271 & 8 & 1 & (8,1) & 10 \\
281 & 283 & 2 & 4 & (2,4) & 22 \\
311 & 313 & 5 & 7 & (5,7) & 15 \\
\bottomrule
\end{tabular}
\caption{DR patterns of the first twenty twin prime pairs.
Only the three patterns $(2,4)$, $(5,7)$, $(8,1)$ appear for $p>3$.}
\label{tab:twin}
\end{table}

%──────────────────────────────────────────────────────────────────────────────
\section{Canonical Degeneracy in the 37-Field}
\label{sec:37field}
%──────────────────────────────────────────────────────────────────────────────

\begin{proposition}[Canonical Degeneracy]
\label{prop:canon}
The three canonical digit blocks satisfy
\[
    123 \equiv 456 \equiv 789 \equiv 12 \pmod{37}.
\]
\end{proposition}

\begin{proof}
$123 = 3\cdot37 + 12$;\quad $456 = 12\cdot37 + 12$;\quad $789 = 21\cdot37 + 12$.
\end{proof}

The common residue 12 is significant: $\mathrm{DR}(12) = 3 \in \mathcal{F}$,
so the 37-field collapses all three canonical blocks onto a value whose digital
root is the first Tesla FLUX state.  Moreover, the digit sum of 12 is $1+2=3$,
and the digits 1 and 2 that appear in 12 are the first two counting numbers,
with $1+2=3$ completing the sequence.

\begin{observation}
$[123]_{37} = [456]_{37} = [789]_{37} = 12 = 1+2+3-\mathrm{DR}(123)$.
\end{observation}

This degeneracy means that under mod-37 arithmetic, the ordered partition of
digits $\{1,\ldots,9\}$ into the blocks $(123,456,789)$ produces no residue
variation --- all three blocks are indistinguishable.  A differential cascade
$\{+8, +13, +24\}$ is required to break this degeneracy, as analysed in the
lattice transform engine described in Appendix~\ref{app:lattice}.

%──────────────────────────────────────────────────────────────────────────────
\section{The 191--137 Prime Bridge}
\label{sec:bridge}
%──────────────────────────────────────────────────────────────────────────────

We record a cluster of arithmetic relations connecting the primes 191 and 137
through the 37-field.

\begin{proposition}[191--137 Bridge]
\label{prop:bridge}
The following identities hold:
\begin{align}
    191 - 137 &= 54 = 6 \times 9, \label{eq:gap}\\
    \mathrm{DR}(54) &= 9, \label{eq:dr54}\\
    54 \bmod 37 &= 17, \label{eq:bridge_seed}\\
    137 + 37 + 17 &= 191, \label{eq:lift}\\
    191 \bmod 37 &= 6 \in \mathcal{F}, \label{eq:tesla}\\
    137 \bmod 37 &= 26,\quad 26+11=37, \label{eq:complement}\\
    \mathrm{DR}(191) &= \mathrm{DR}(137) = 2. \label{eq:resonance}
\end{align}
Moreover, 191 is the 43rd prime and 137 is the 33rd prime, with
$43 - 33 = 10$ and $43\bmod 37 = 6 = [191]_{37}$.
\end{proposition}

\begin{proof}
All identities are direct computation.
\end{proof}

\begin{observation}[Twin prime pairs at both bridge seeds]
Both bridge seeds lead twin prime pairs:
$(137, 139)$ and $(191, 193)$ are each twin prime pairs.
By Corollary~\ref{cor:twin_dr2}, since $\mathrm{DR}(137)=\mathrm{DR}(191)=2$,
both pairs carry pattern $(2,4)$.  The shared DR resonance of the bridge
is exactly the leading DR value of both twin prime pairs.
\end{observation}

In the 37-field the twin prime gap of 2 is preserved:
$[191]_{37}=6$ and $[193]_{37}=8$, so $8-6=2$;
$[137]_{37}=26$ and $[139]_{37}=28$, so $28-26=2$.

%──────────────────────────────────────────────────────────────────────────────
\section{Pascal's Row 8: Unified Analysis}
\label{sec:pascal}
%──────────────────────────────────────────────────────────────────────────────

Pascal's Row 8 is
\[
    R_8 = (1,\;8,\;28,\;56,\;70,\;56,\;28,\;8,\;1),
    \qquad \sum R_8 = 256 = 2^8.
\]
Its element digital roots are $\mathrm{DR}(R_8) = (1,8,1,2,7,2,1,8,1)$, giving
spine FLUX ratio 0 --- no element is FLUX-active.

\subsection{The 37-Field Pivot}

\begin{proposition}
\label{prop:pascal37}
The first three elements of $R_8$ sum to the 37-field modulus:
\[
    1 + 8 + 28 = 37.
\]
The running prefix sum hits $[s]_{37}=19$ at position 3
(after including $r_3=56$): $1+8+28+56=93$ and $93\bmod 37 = 19$.
\end{proposition}

\begin{proof}
$1+8+28=37$ and $93=2\cdot37+19$.
\end{proof}

The value 19 is significant: $\mathrm{DR}(19)=1$ and the digits of 19 are
$\{1,9\}$, matching the digit pattern of the prime 191.

\subsection{The 191 Bridge Intersection}

\begin{proposition}
\label{prop:pascal191}
In the 37-field, the elements of $R_8$ have residues
\[
    (R_8 \bmod 37) = (1,\;8,\;28,\;19,\;33,\;19,\;28,\;8,\;1).
\]
Specifically:
\begin{itemize}
\item $[56]_{37} = 19$ (positions 3 and 5): the 191 digit pattern $\{1,9\}$
      appears as a 37-field residue.
\item $[70]_{37} = 33$ and $\mathrm{DR}(33) = 6 = [191]_{37}$: the centre
      element maps to the same Tesla FLUX class as the prime 191.
\end{itemize}
\end{proposition}

\begin{proof}
$56 = 37+19$;\quad $70 = 37+33$;\quad $191 = 5\cdot37+6$;\quad $\mathrm{DR}(33)=6$.
\end{proof}

\subsection{Tesla FLUX Saturation under the Pascal Transition}

\begin{theorem}[Adjacent-pair FLUX saturation]
\label{thm:pascal_flux}
Every adjacent pair sum in $R_8$ has digital root in $\mathcal{F}$:
\[
    \mathrm{DR}(r_i + r_{i+1}) \in \{3,6,9\}
    \quad \text{for } i = 0,1,\ldots,7.
\]
The spine FLUX ratio of the adjacent pair sums is $8/8 = 1$.
\end{theorem}

\begin{proof}
The adjacent pair sums are
\[
    (9,\;36,\;84,\;126,\;126,\;84,\;36,\;9)
\]
with digital roots $(9,9,3,9,9,3,9,9)$.  All lie in $\mathcal{F}$.
\end{proof}

The adjacent pair sums of $R_8$ are exactly the interior elements of $R_9$
(Pascal's Row 9).  Theorem~\ref{thm:pascal_flux} states that the Pascal
transition $R_8 \to R_9$ maps every interior position to a FLUX-active value:
the operation by which Pascal's triangle builds the next row saturates the
Tesla spine at 100\%.

\subsection{The Accumulation Chain}

Running the cumulative sum $S_k = \sum_{i=0}^{k} r_i$ and computing
$\mathrm{DR}(S_k)$:

\begin{center}
\begin{tabular}{rrrcc}
\toprule
$k$ & $r_k$ & $S_k$ & $[S_k]_{37}$ & $\mathrm{DR}(S_k)$ \\
\midrule
0 & 1  &   1 &  1 & 1 \\
1 & 8  &   9 &  9 & 9 (\textbf{FLUX}) \\
2 & 28 &  37 &  0 & 1 (37-field zero) \\
3 & 56 &  93 & 19 & 3 (\textbf{FLUX}) \\
4 & 70 & 163 & 15 & 1 \\
5 & 56 & 219 & 34 & 3 (\textbf{FLUX}) \\
6 & 28 & 247 & 25 & 4 \\
7 & 8  & 255 & 33 & 3 (\textbf{FLUX}) \\
8 & 1  & 256 & 34 & 4 \\
\bottomrule
\end{tabular}
\end{center}

The prefix sum passes through the 37-field zero at $k=2$ ($S_2=37$),
touches the 191-pattern residue 19 at $k=3$ ($[S_3]_{37}=19$), and
achieves FLUX-active DR at positions $k\in\{1,3,5,7\}$ --- the four
odd positions.  At position 7, $[S_7]_{37}=33$ and $\mathrm{DR}(33)=6$,
the Tesla class shared with $[191]_{37}$.

%──────────────────────────────────────────────────────────────────────────────
\section{The Digit-Sum Invariant and the 123 Fixed Point}
\label{sec:digitsum}
%──────────────────────────────────────────────────────────────────────────────

\begin{theorem}[Digit-sum invariant]
\label{thm:digitsum}
Let $a, b \geq 1$ with $\sigma(a) + \sigma(b) = 12$, where $\sigma(n)$
denotes the digit sum of $n$.  Then $\sigma(a+b) \equiv 12 \pmod{9}$.
If $a+b < 100$, then $\sigma(a+b) = 12$.
\end{theorem}

\begin{proof}
Since $\sigma(n) \equiv n \pmod{9}$, we have $\sigma(a+b) \equiv a+b
\equiv \sigma(a)+\sigma(b) = 12 \equiv 3 \pmod{9}$.  For $a+b < 100$,
the digit sum is between 1 and 18, and the unique value in that range
that is $\equiv 3\pmod 9$ is $12$ (since $3<12<18$ and neither 3 nor
21 lies in $[1,18]$ with the right residue --- 3 and 12 both satisfy,
but $a+b\geq 20$ whenever both $\sigma(a)\geq 2$ and $\sigma(b)\geq 2$
with their sums totalling 12, so 3 does not arise for $a+b<100$).
\end{proof}

\begin{observation}
$12 = [123]_{37} = [456]_{37} = [789]_{37}$.
\end{observation}

The fixed point 12 is precisely the canonical mod-37 degeneracy of
Section~\ref{sec:37field}.  The three examples below use the digit set
$\{1,3,4,4\}$ (digit sum 12) rearranged into valid additions:
\[
    13+44=57,\quad 144+3=147,\quad 134+4=138.
\]
In each case $\sigma(\text{left side}) = 12$ and $\sigma(\text{right side}) = 12$.
The 123 residue is invariant under rearrangement of the digit set.

%──────────────────────────────────────────────────────────────────────────────
\section{DR-Product Chain for Row 8}
\label{sec:drproduct}
%──────────────────────────────────────────────────────────────────────────────

\begin{theorem}[DR-product chain]
\label{thm:drproduct}
Let $P = \prod_{i=0}^{8} \mathrm{DR}(r_i)$ be the product of the digital
roots of $R_8$.  Then
\[
    P = 1792,\qquad \sigma(P) = 19,\qquad \mathrm{DR}(P) = 1.
\]
Moreover $P = 2^8 \cdot \mathrm{DR}(r_4)$, where $r_4 = 70$ is the centre
element.
\end{theorem}

\begin{proof}
$\mathrm{DR}(R_8) = (1,8,1,2,7,2,1,8,1)$.  The non-unit factors are
$8,2,7,2,8$, giving $P = 8\cdot2\cdot7\cdot2\cdot8 = 1792$.
$\sigma(1792)=1+7+9+2=19$.  $\mathrm{DR}(1792)=\mathrm{DR}(19)=1+9=10
\to 1$.  Finally $1792 = 256\cdot7 = 2^8\cdot7 = 2^8\cdot\mathrm{DR}(70)$.
\end{proof}

The chain $1792 \to 19 \to 1$ passes through the intermediate value 19,
whose digits $\{1,9\}$ match the digit pattern of the prime seed 191.
The same pattern 19 appears as the 37-field residue of the elements at
positions 3 and 5 of $R_8$ (Proposition~\ref{prop:pascal191}), and as the
prefix-sum residue at position 3 of the accumulation chain
(Proposition~\ref{prop:pascal37}).

\begin{corollary}
The DR-product chain for $R_8$ terminates at 1.  The intermediate value 19
is the 191-digit pattern appearing in the 37-field landscape of $R_8$.
\end{corollary}

%──────────────────────────────────────────────────────────────────────────────
\section{The Twin Prime Pair $(41,43)$ and the 123 Chain}
\label{sec:4143}
%──────────────────────────────────────────────────────────────────────────────

The twin prime pair $(41,43)$ is distinguished within this framework by
two facts: (a) 43 is the prime index of 191 ($P(43)=191$), and (b)
$43\bmod 37 = 6 = [191]_{37}$, so the index shares the Tesla class of
the seed.

\begin{proposition}
$41+43=84$ and the digit-sum chain of 84 passes through
$1,2,3$ sequentially:
\[
    84 \;\xrightarrow{\sigma}\; 12 \;\xrightarrow{\sigma}\; 3.
\]
The intermediate value 12 equals $[123]_{37}$, and the terminal value
3 is the first Tesla FLUX state.
\end{proposition}

\begin{proof}
$41+43=84$;\; $\sigma(84)=8+4=12=[123]_{37}$;\; $\sigma(12)=1+2=3\in\mathcal{F}$.
\end{proof}

The digits visited in the chain are $8,4,1,2,3$ --- the digit 3 is
reached only at the end, as the FLUX attractor.  The digit sequence
$1,2,3$ appears explicitly in the reduction.

%──────────────────────────────────────────────────────────────────────────────
\section{The Zero Principle and the Expanding Pyramid}
\label{sec:zero}
%──────────────────────────────────────────────────────────────────────────────

Pascal's triangle has a natural zero state: Row 0 is the singleton $[1]$,
representing the first knowing --- existence before categorisation.  The
subsequent rows introduce the counting numbers in order:

\begin{center}
\begin{tabular}{cl}
\toprule
Row & First new digit sum value \\
\midrule
0 & 1 (the unit) \\
2 & 2 (first occurs in $\binom{2}{1}=2$) \\
3 & 3 (first occurs in $\binom{3}{1}=3$) \\
4 & 6 $\in\mathcal{F}$ (first occurs in $\binom{4}{2}=6$) \\
\bottomrule
\end{tabular}
\end{center}

The sequence $1\to 2\to 3\to 6$ traces the pyramid's structural unfolding:
the counting numbers 1, 2, 3 precede the first Tesla FLUX state 6, which
first appears at Row 4.  Since $1+2+3=6$ and $1\cdot2\cdot3=6$, the
appearance of 6 at Row 4 is the first moment the cumulative and
multiplicative structures of $\{1,2,3\}$ are realised as a binomial coefficient.

%──────────────────────────────────────────────────────────────────────────────
\section{Conclusion and Open Questions}
\label{sec:conclusion}
%──────────────────────────────────────────────────────────────────────────────

We have established the following results:

\begin{enumerate}
\item (\S\ref{sec:twin}) The only DR patterns for twin prime pairs $(p,p+2)$
with $p>3$ are $(2,4)$, $(5,7)$, $(8,1)$.  The Tesla FLUX class $\{3,6,9\}$
is structurally excluded.

\item (\S\ref{sec:37field}) The canonical digit blocks 123, 456, 789 all
reduce to 12 mod 37.  The value 12 is the fixed point linking the 123
digit pattern to the 37-field.

\item (\S\ref{sec:bridge}) The primes 191 and 137 share DR 2, satisfy
$191-137=6\times 9$, and each leads a twin prime pair with pattern $(2,4)$.
The prime index of 191 is 43, and $43\bmod 37 = 6 = [191]_{37}$.

\item (\S\ref{sec:pascal}) Pascal's Row 8 satisfies $1+8+28=37$ (37-field
pivot), has centre element with $[70]_{37}=33$ and $\mathrm{DR}(33)=6=[191]_{37}$,
achieves 100\% Tesla FLUX in adjacent pair sums, and its DR-product chain
resolves through 19 to 1.

\item (\S\ref{sec:digitsum}) Any digit set summing to 12 generates additions
whose sums also have digit sum 12.  The fixed point $12=[123]_{37}$ is
preserved by the digit-sum operation.

\item (\S\ref{sec:4143}) The sum of the twin prime pair $(41,43)$ reduces
via 12 (the 123-residue) to 3 (the first Tesla FLUX state).
\end{enumerate}

\subsection*{Open questions}

\begin{enumerate}
\item \textbf{Twin prime distribution across patterns.}  Numerical evidence
suggests the three DR patterns $(2,4)$, $(5,7)$, $(8,1)$ occur with
roughly equal frequency.  Does each pattern contain infinitely many twin
prime pairs (conditional on the twin prime conjecture)?

\item \textbf{37-field structure.}  What is the distribution of twin prime
residues mod 37?  Are there 37-field residue classes that are avoided?

\item \textbf{Pascal rows and FLUX saturation.}  Theorem~\ref{thm:pascal_flux}
shows that Row 8's adjacent pair sums achieve 100\% FLUX.  For which rows $n$
do the adjacent pair sums achieve full FLUX saturation?  Is there a
characterisation in terms of $n\bmod 9$?

\item \textbf{Generalisation of the DR-product chain.}  For which rows of
Pascal's triangle does the product of element digital roots reduce to 1?
\end{enumerate}

%──────────────────────────────────────────────────────────────────────────────
\appendix
\section{Lattice Transform Engine}
\label{app:lattice}
%──────────────────────────────────────────────────────────────────────────────

The 1,680 ordered 3-partitions of $\{1,\ldots,9\}$ into groups of three
are generated by $\binom{9}{3}\binom{6}{3}\binom{3}{3} = 84\cdot20\cdot1$.
Each partition $(g_1,g_2,g_3)$ defines block values $b_i$ (the three-digit
number formed by the sorted digits of $g_i$).  The differential cascade
$b_i \mapsto b_i + c_i$ with $(c_1,c_2,c_3) = (8,13,24)$ breaks the
mod-37 degeneracy of Proposition~\ref{prop:canon}: computational search
finds that 35 of the 1,680 partitions survive the adjacency filter requiring
all pairwise differences of cascaded residues to lie in $\{8,13,24\}$ or
their mod-37 complements $\{29,24,13\}$.  The cascade constants satisfy
$8+13+24 = 45 = \sum_{k=1}^{9} k$ and $\mathrm{DR}(8)+\mathrm{DR}(13)+\mathrm{DR}(24)
= 8+4+6 = 18$, with $\mathrm{DR}(18)=9\in\mathcal{F}$.

\end{document}
